\chapter*{Preface}
\addcontentsline{toc}{chapter}{Preface}

Meet Alex. It is Sunday evening, and Alex's desk is covered in a messy pile of homework sheets, a notebook tracking weekly pocket money, and a sticky note reminder to wake up early for school. Alex looks at the mess and thinks, \textit{"If only I could tell someone to organize this for me automatically..."}

That is what programming is all about.

Programming is not just for computer scientists or people in movies. Today, programming is a basic life skill—just like reading, writing, and mathematics. Every day, you use mobile phones, apps, and computers, but very few people understand how they actually work under the hood.

This book is written to help you:
\begin{itemize}
    \item Think clearly and logically.
    \item Solve everyday problems using code.
    \item Build confidence as a beginner.
\end{itemize}

This textbook is specially designed for high school students with no prior coding experience, following modern school curriculums.

\section*{What Makes This Book Different?}

Most programming books focus only on strict rules and hard words. This book focuses on something different:
\begin{itemize}
    \item \textbf{Real-life stories:} Every chapter starts with a situation you face every day.
    \item \textbf{Storytelling thread:} You will follow characters like Alex and Maya as they solve daily puzzles.
    \item \textbf{Step-by-step learning:} Concepts grow naturally from the stories, one small piece at a time.
\end{itemize}

Programming is a \textbf{tool for thinking}, not just typing code.

\section*{Why Python?}

In this book, we use **Python**. Python is a programming language that looks almost like normal English. It is easy to read, simple to write, and powerful enough to build real applications.

By using Python, you can focus entirely on \textbf{your logic} instead of getting stuck on difficult symbols.

\section*{What You Will Achieve}

By the time you finish this book, you will be able to:
\begin{itemize}
    \item Understand how computer programs make decisions.
    \item Write your own small, useful programs.
    \item Think systematically when facing a hard problem.
\end{itemize}

This book does not aim to turn you into a software engineer overnight. It aims to make you a \textbf{confident beginner} who is ready for the digital world.

\section*{A Quick Note}

Do not fear making mistakes. In programming, errors are just clues that help you learn.

Read slowly, practice every day, and enjoy the journey!

\vfill
\begin{flushright}
\textit{The Author}
\end{flushright}

\newpage

\chapter*{How to Use This Book}
\addcontentsline{toc}{chapter}{How to Use This Book}

This book is designed for both classroom learning and self-study at home. Follow the steps below to get the best results.

\section*{For Students: Your Learning Path}

To help you navigate each chapter smoothly, follow this visual workflow:

\begin{center}
\begin{tikzpicture}[node distance=1.5cm, auto]
    \node (story) [startstop, text width=5cm] {1. Read the Story\\(Meet Alex \& Maya)};
    \node (concept) [process, text width=5cm, below of=story] {2. Understand the Concept\\(Why and where it is used)};
    \node (example) [process, text width=5cm, below of=concept] {3. Study the Code Examples\\(Type them yourself)};
    \node (exercise) [process, text width=5cm, below of=example] {4. Practice Exercises\\(Build muscle memory)};
    \node (project) [startstop, text width=5cm, below of=exercise] {5. Try the Mini-Project\\(Solve a real-life task)};

    \draw [arrow] (story) -- (concept);
    \draw [arrow] (concept) -- (example);
    \draw [arrow] (example) -- (exercise);
    \draw [arrow] (exercise) -- (project);
\end{tikzpicture}
\end{center}

\subsection*{A Few Golden Rules}
\begin{itemize}
    \item \textbf{Do not skip stories:} They explain \textit{why} you need the code.
    \item \textbf{Type the code yourself:} Never just read code like a novel. Typing builds your coding muscle memory.
    \item \textbf{Embrace errors:} When a red error message pops up on your screen, don't panic. It is just your computer asking you to check your logic.
\end{itemize}

\section*{For Teachers}

This textbook supports active, discussion-based learning. Teachers are encouraged to:
\begin{itemize}
    \item Connect every topic back to the chapter's opening story.
    \item Let students experiment and break code to see what happens.
    \item Focus heavily on logic and problem-solving over memorizing syntax.
\end{itemize}

\section*{Symbols Used in This Book}

\begin{itemize}
    \item \textbf{Example} — Displays sample Python code and output.
    \item \textbf{Mini-Project} — A complete, real-life application built step-by-step.
    \item \textbf{Activity} — Quick practice tasks for homework or classwork.
\end{itemize}

Take your time, think clearly, and happy coding!

\vfill
\begin{flushright}
\textit{The Author}
\end{flushright}